Can a deterministic multitape Turing machine transpose a $d \times d$ matrix with $\max(1, \clog{d})$-bit entries in time
linear in the input length? By the fourth paper of this series, excluding oblivious linear-time transposers splits into two
combinatorial problems: a half-rate bound for storage codes on transitive-tournament-free recovery relations (the coding
crux), and cheap depth reduction of head-walk dependency graphs (the trajectory cut problem). This paper advances three
frontiers left open by the seventh, eighth and ninth papers of the series, and every new numbered theorem, proposition,
lemma and corollary is kernel-verified in Lean~4.

First, for nested elimination on directed recovery relations, we give a purely positional characterization of the
nestedness hypothesis of the seventh paper, show that every reciprocal pair forces distinct record lists, and prove a
sharp dichotomy for chordless digon cycles with at most one active outside parent per vertex: a piece certificate exists
over some number of records if and only if the outside arcs do not form a unidirectional cycle of stride~$3$
(Theorem~DC). In particular, the $15$-vertex obstruction $\Ex$ has no piece certificate on its unique minimum nonempty
core over any number of records, and the Petersen graph $\GP(5,2)$ has no global elimination certificate over any number
of records, whereas $\GP(7,2)$ and $\GP(9,2)$ require minimum record counts $4$ and~$3$, respectively (Proposition~GP),
resolving two open formalization problems of the seventh paper. Using minimal chordal digon paths and hub-augmented
stride-$3$ pieces, we then prove that the entire class $\ClassS$ of Theorem~S of the fourth paper---transitive-tournament-free
relations of maximum in-degree~$3$ in which every vertex of in-degree~$3$ has only reciprocal children---admits a
certified piece decomposition and satisfies the half-rate bound $\card{C}^2 \le \card{\alpha}^{\card{V}}$ over every
finite alphabet and for all codes, linear or nonlinear (Theorem~S$^*$), resolving the all-code Theorem~S target of the
seventh paper, together with a directed barrier inequality (Corollary~B$^*$) and a $21$-vertex strongly connected
triangle-free witness $\Gto \in \ClassS$ on which earlier peeling rules stall (Proposition~G21).

Second, for scalar linear codes on triangle-free graphs of maximum degree~$4$ defined as the kernel of a matrix $M$ with
unit diagonal over an arbitrary field, we remove the heavy-separation hypothesis~(S2) of the eighth paper's Theorem~SH:
adjacent degree-$4$ exceptional vertices and arbitrary matchings of exceptional vertices are allowed, provided no two
exceptional vertices share a common neighbor (Theorem~SH$^*$), yielding a matching bound controlled by the heavy-matching
invariant $\rhom(G,H) \ge \rho_3(G,H)$ (Corollary~M$^*$).

Third, on the deep-nesting walk $U(s_0, m, 1)$, we formalize the recursive stroke structure and $O(m)$ turn-potential
evaluator (Proposition~SL), prove an $m$-free mod-$q$ straddle lower bound and exact parity-conditioned straddle count
$\Pi_\kappa^2 / 4$ (Theorem~V$^*$), establish a super-Fibonacci unwrapping theorem showing that level-uniform single-head
cancellation admits at most $2\lfloor \log_2 q \rfloor + 2$ consecutive levels whose blind apex and band-$0,1,2$ errors are
at most $q/8$ (Theorem~LU), determine the block-cancellation drift and Mersenne $q = 2^B - 1$ binary-clock failure residue
(Proposition~OC), and prove an unconditional three-regime witness theorem with zero $\log q$ loss for block-label pairs of
bounded block length across every $K_1, K_2 \in \Z$ and every modulus $q \ge 3$ at two fixed odd shifts, or a single shift
for symmetric pairs (Theorem~A; translating its single-head regimes to actual two-head hop errors requires blindness at the
chosen shift). The linear-time transposition question itself remains open.
